% Chapter 8 -- The InitechDOS Editing Keys.
% Law 1: MS-DOS 3.3 User's Guide printed pp.153-159,165-169
% (PDF238-244,250-254); ADR-0003 Appendix A (01h-0Ch CON I/O).
% conedit.h is shared by int21.c AH=0Ah/CON and EDLIN; command.c retains
% the previous buffered line. Printer hardware is outside spec/hardware.json:
% PRN uses the existing serial backend, not an invented parallel port.
% GATES: test-conkeys-unit test-conkeys test-conkeys-mutant test-edlin-emu
\Gates{test-conkeys-unit, test-conkeys, test-conkeys-mutant, test-edlin-emu}
\chapter{The InitechDOS Editing Keys}

\section{Typing a Command Line}
\index{editing keys}
InitechDOS reads what you type one line at a time. The characters appear on
 the screen as you type them, and nothing happens until you press
\key{Enter}. Until then you can correct the line with the keys described
in this chapter. The same editing keys work when a program asks you to
type a line, and when you edit a line with \cmd{EDLIN}.

At the command prompt, a line can be up to 127 characters long. When it is
full, InitechDOS sounds a beep and ignores further characters until you
press \key{Enter} or erase some.

\section{The Keys}
\SideHead{Enter}Ends the line. InitechDOS moves the cursor to the next
line and carries out the command. The completed line becomes the template
for the next command.

\SideHead{Backspace}Erases the character to the left of the cursor and
moves back one character in the template when you are replacing characters.
At the beginning of the line it does nothing. \key{Ctrl-H} does the same.

\SideHead{Ctrl-C}Cancels the line you are typing: InitechDOS discards it
and displays the command prompt again. If a command or program is waiting
for you to type a line, \key{Ctrl-C} stops it.  In a batch job, InitechDOS
asks \texttt{Terminate batch job (Y/N)?}: type \type{Y} to end the job
or \type{N} to continue with the next command.
InitechDOS displays \texttt{\textasciicircum C} after the cancelled line:
\begin{Steps}
\item Type \type{CLS} and press \key{Enter}.
\item Type \type{DIR}, but do not press \key{Enter}: press \key{Ctrl-C}
  instead.
\end{Steps}

\Screen[rows=2, label=fig:ctrlc]{screens/ctrlc.txt}{A line cancelled with
  Ctrl-C}

\SideHead{Ctrl-Z}Marks the end of a file. When you copy from the keyboard
with \cmd{COPY CON}, press \key{Ctrl-Z} and then \key{Enter} to end the
file.

\section{Correcting a Line with Backspace}
To correct a mistake, press \key{Backspace} until the wrong character is
erased, then type the rest of the line again. For example, if you type
\type{DIT}, press \key{Backspace} once and type \type{R}: the line becomes
\type{DIR}, and pressing \key{Enter} lists the directory.

\section{The Template Keys}
InitechDOS keeps the last completed line as a \term{template}. A position
in the template marks the next character available to copy. Normally, each
character you type replaces one character at that position.

\SideHead{F1}Copies one character from the template to the current line.

\SideHead{F2}Copies characters up to, but not including, the character you
press next. If that character is absent, nothing is copied.

\SideHead{F3}Copies the rest of the template. Press \key{F3} and
\key{Enter} at the prompt to repeat the last command. \key{F1}, \key{F2}
and \key{F3} also turn off insert mode.

\SideHead{F4}Skips characters in the template up to, but not including,
the character you press next. If it is absent, nothing is skipped.

\SideHead{F5}Makes the current line the new template without carrying it
out. InitechDOS prints \texttt{@} and moves to a new line, ready for you to
copy or change this template. Insert mode is turned off.

\SideHead{Del}Skips one character in the template without copying it to
the current line. The cursor on the screen does not move.

\SideHead{Ins}Turns insert mode on or off. In insert mode, characters you
type are added to the current line without moving through the template.
Press \key{Ins} again to return to replacing characters.

\SideHead{Esc}Cancels the current line without changing the template.
InitechDOS prints a backslash (\texttt{\textbackslash}) and moves to a new
line. Insert mode is turned off. Press \key{F3} to copy the template again.

Try the following example:
\begin{Steps}
\item Type \type{CLS}, then \type{ECHO Quarter}, pressing \key{Enter}
  after each command.
\item Press \key{F3} and \key{Enter} to repeat the command.
\item Press \key{F2}, then \key{Q}. Type \type{Month} and press \key{Enter}.
\item Type \type{ECHO Friday} and press \key{F5}. Press \key{F3} and
  \key{Enter} to carry out the saved line.
\item Type \type{WRONG} and press \key{Esc}. Press \key{F3} and
  \key{Enter} to recall the unchanged template.
\end{Steps}

\Screen[rows=13, label=fig:template]{screens/template.txt}{Copying and
  cancelling lines with the template keys}

\section{Controlling the Display and Printer Echo}
\SideHead{Ctrl-S}Pauses the display. Press any key to resume; that key is
used only to resume the display.

\SideHead{Ctrl-P}Turns printer echo on or off. While it is on, output sent
to the console is also sent to \texttt{PRN}. This edition records printer
output on the diagnostic serial connection; physical printing requires a
printer driver.

\SideHead{Ctrl-N}Turns printer echo off.

